% !PRINTED none

\begin{center}
\gesperrt{Eleventh Section.}
\end{center}

\begin{center}
\gesperrt{Continued Fractions.}
\end{center}

\begin{center}
\rule{3cm}{0.4pt}
\end{center}

\begin{center}
§. 115.
\end{center}

\begin{center}
Transformation of rational fractions into continued fractions.
\end{center}

If one divides a natural number $m$ by another $n$ according to the ordinary rules, one obtains a quotient and a remainder, which is positive and smaller than the divisor $n$. If $a$ denotes the quotient and $r$ the remainder, then
\weq{1}{m = an + r.}
All numbers $m, m', m''\ldots$ which give the same remainder are called \gesperrt{of equal residue} or \gesperrt{congruent} with respect to the \gesperrt{modulus} $n$, and following \gesperrt{Gauss} one expresses this by the sign:
\weq{2}{m \equiv m' \ (\mathrm{mod}\ n).}
Such a formula is called a \gesperrt{congruence}.

If, then, we consider the system of numbers
\[
r = 0, 1, 2 \ldots n - 1
\]
every arbitrary (positive or negative) integer $m$ is congruent to \gesperrt{one} and only \gesperrt{one} of these numbers with respect to the modulus $n$. But the same property is also possessed by every system of numbers
\[
s = m_0, m_1, m_2 \ldots m_{n-1},
\]
chosen in such a way that among their remainders all the numbers $r$ occur (and each only once). Then every arbitrary number $m$ is, with respect to the modulus $n$, congruent to one and only one of the numbers $s$
%CONT>>
